\fichatitulo{23 · RECUPERAÇÃO}{Preprint v1 · 2020}{ColBERT}
{\small Khattab; Zaharia. \textit{{ColBERT}: Efficient and Effective Passage Search via Contextualized Late Interaction over {BERT}}. Preprint v1 · 2020. \href{https://arxiv.org/pdf/2004.12832v1}{Fonte consultada}.\par}

\campo{Problema e objetivo}
Como combinar interação contextual e eficiência na busca?

\campo{Principal contribuição · síntese interpretativa}
Propõe interação tardia entre representações de tokens, permitindo pré-computar as representações dos documentos.

\campo{Método / natureza da evidência}
Codifica consulta e documento separadamente; calcula interações entre tokens para pontuar a relevância.

\campo{Resultado ou argumento principal}
O texto apresenta uma alternativa entre representações únicas por documento e modelos que processam conjuntamente cada par consulta-documento.

\campo{Excertos-chave · seleção literal}
\begin{itemize}
\excerto{1}{pre-compute document representations offline}{Resumo.}
\excerto{2}{late interaction architecture}{Resumo.}
\end{itemize}

\campo{Limitações e análise crítica}
Análise da ficha: múltiplos vetores por documento implicam custos de armazenamento e busca. Comparações devem explicitar índice, hardware e qualidade, sem inferir superioridade para qualquer corpus.

\campo{Aproveitamento na dissertação · proposta}
Fundamentação e metodologia: justificar uma eventual comparação entre recuperação lexical, densa e interação tardia.

\registro{Fonte consultada nesta preparação: Resumo e introdução. Chave: \texttt{khattabZaharia2020colbert}.}
